\section{Methods}
\label{sec:methods}

\subsection{DFT data, equations of state, and formation energies}

The repository stores a curated Fe--Mo dataset assembled from VASP calculations using the projector-augmented-wave method and the Perdew--Burke--Ernzerhof exchange-correlation functional. The nonmagnetic workflow archived in \texttt{Fe-Mo/inchulldft/setup.yaml} is labeled \texttt{VASP\_PBE\_500\_0.125\_0.1\_NM} and specifies \texttt{ENCUT}=\SI{500}{\electronvolt}, a $\Gamma$-centered reciprocal-space spacing of \SI{0.125}{\per\angstrom}, \texttt{ISPIN}=0, \texttt{ISMEAR}=0, \texttt{SIGMA}=\SI{0.2}{\electronvolt}, and tight electronic convergence with \texttt{EDIFF}=$10^{-6}$. The curated notebook output further shows that the final retained structures are dominated by closely related PAW--PBE relaxations with \texttt{ENCUT}=\SI{450}{\electronvolt} and an effective $\Delta k\approx\SI{0.02}{\per\angstrom}$, while a small number of elemental reference structures were retained at neighboring cutoffs. The relaxation metadata and the curation notebook both identify Murnaghan fits as the standard equation-of-state treatment used to obtain relaxed volumes and bulk moduli.\cite{KresseHafner1993,KresseFurthmueller1996CMS,KresseFurthmueller1996PRB,Perdew1996,Bloechl1994,KresseJoubert1999}

For each structure, the notebooks tabulate formation energies relative to several magnetic reference choices. The target used throughout the machine-learning workflow is $E_f^{\mathrm{nmhcp}}$, namely the formation energy with respect to nonmagnetic hcp Fe and hcp Mo. With $N_{\mathrm{Fe}}$ Fe atoms and $N_{\mathrm{Mo}}$ Mo atoms in a structure containing $N=N_{\mathrm{Fe}}+N_{\mathrm{Mo}}$ atoms, the definition is
\begin{equation}
E_f^{\mathrm{nmhcp}} = \frac{E_{\mathrm{tot}} - N_{\mathrm{Fe}}\mu_{\mathrm{Fe}}^{\mathrm{nmhcp}} - N_{\mathrm{Mo}}\mu_{\mathrm{Mo}}^{\mathrm{hcp}}}{N}.
\label{eq:formation_energy}
\end{equation}
The final curation table reports 292 usable structures after one training-set removal, with $N$ ranging from 1 to 53 atoms and with $E_f^{\mathrm{nmhcp}}$ spanning from \SI{-42.841}{\milli\electronvolt\per\atom} to \SI{720.492}{\milli\electronvolt\per\atom}. The Zenodo manifest bundled with the repository identifies the corresponding curated DFT archive as containing 291 training structures, which resolves the apparent 292/291 ambiguity by making explicit that one curated entry is retained in the notebook summary but excluded from the training set. The downstream regression analysis excludes bcc, fcc, and hcp reference phases and evaluates 261 remaining structures, while the validation notebook adds 70 complex-phase calculations comprising 19 $R$, 12 $P$, 11 $M$, and 28 $\delta$ structures. The same manifest separates the raw complex-phase validation inputs into 14 entries under \texttt{Fe-Mo/validation\_data} and 21 further in-hull entries under \texttt{Fe-Mo/inchulldft}, which are merged by the validation notebook before parity analysis.

The simple TCP training manifold is built around A15, C15, C14, C36, $\sigma$, $\chi$, and $\mu$ prototypes, whereas the complex-phase validation set targets $R$, $M$, $P$, and $\delta$. The coordination dictionaries shipped with the repository make the phase complexity explicit. For example, C14 contains three coordination-site families with multiplicities $[4,2,6]$, $\mu$ contains five with multiplicities $[1,6,2,2,2]$, and $\sigma$ contains five with multiplicities $[2,4,8,8,8]$, while the $R$ phase requires eleven sublattices with multiplicities $[1,2,6,6,6,6,6,6,6,2,6]$. Initial structures for compositional sweeps are generated by Vegard-type volume scaling before relaxation, so that the subsequent DFT optimization begins from chemically sensible cell volumes rather than from fixed end-member lattices.

\subsection{Descriptor construction and coordination-number-resolved averaging}

Three classes of atom-centered descriptors were compared. The most specialized representation uses bond-order-potential moments derived from a canonical $d$-band tight-binding description, which connects the moments of the local density of states to the geometry of the surrounding bonding network. In the archived feature files the best-performing variant is the bond-specific CNAV table \texttt{0.7dProjections 0.5OS BOP}, while a less specialized \texttt{Canonical BOP} representation serves as a reference. This distinction is important because the bond-specific model keeps more of the local-orbital information that is expected to control the stability of transition-metal TCP phases.\cite{PettiforAoki1996,Pettifor2004}

As complementary machine-learning descriptors, the repository includes ACE and SOAP features. The archive names the corresponding CNAV files \texttt{Fe-Mo-ACE-CNAV.csv} and \texttt{CNAV\_soap\_features\_\_specific\_\_r\_cut\_4\_\_n\_max\_5\_\_l\_max\_4\_\_sigma\_0.1\_\_rbf\_gto\_\_periodic\_True.csv}. The source code also contains a SOAP constructor with $r_{\mathrm{cut}}=\SI{4}{\angstrom}$, $n_{\max}=10$, $l_{\max}=5$, $\sigma=0.1$, a Gaussian-type radial basis, and periodic boundary conditions.\cite{Drautz2019,Bartok2013,Himanen2020,Lysogorskiy2021}

The central structural idea is coordination-number-resolved averaging. Instead of averaging an atom-centered descriptor over the full unit cell, the featurizer first groups atoms by their coordination families, such as CN12, CN13, CN14, CN15, and CN16. Descriptor components are then averaged within each family and concatenated into a single structure vector. In TCP phases this procedure preserves information about which crystallographic sublattice contributes a given local environment, and it therefore retains much of the order information that would be destroyed by conventional unit-cell averaging. The repository automatically constructs comparison features with and without this grouping, which makes it possible to isolate the contribution of crystallographic site resolution to the final error.

\subsection{Regression models, feature selection, and data splitting}

The learning workflow is implemented with scikit-learn pipelines.\cite{Pedregosa2011} The model-selection notebook constructs pipelines for kernel ridge regression, multilayer perceptrons, random forests, and Gaussian processes. The archived JSON option files show that kernel ridge regression was optimized over polynomial and, in the magnetic/nonmagnetic variants, also radial-basis kernels; the best fitted BOP model displayed in the stored estimator representation uses \texttt{KernelRidge(}\allowbreak\texttt{alpha=0.1, degree=4, kernel='polynomial')} while a SOAP example uses \texttt{KernelRidge(}\allowbreak\texttt{alpha=0.1, coef0=5, kernel='polynomial')}. The random-forest options vary the tree depth and leaf constraints, and the MLP options use a standardized input representation together with logistic activations, \texttt{lbfgs} optimization, and hidden-layer candidates of $(20,4)$ and $(40,4)$. The ensemble stage repeats the recursive-selection workflow ten times (\texttt{n\_repeats = 10} in the archived driver), stores the resulting regressors in an indexed bag, and passes those ten fitted pipelines to a \texttt{VotingRegressor}. The validation notebooks expose the corresponding vote spread through the column \texttt{std\_votes}, which we interpret as an internal measure of ensemble dispersion rather than as a calibrated statistical confidence interval.

Forward recursive feature selection is nested inside the training workflow. The feature-concatenation driver applies a phase-stratified 20\% hold-out split with random state 42 and then performs a fivefold \texttt{StratifiedKFold} cross-validation on the training subset while ranking candidate feature additions by the negative root-mean-square error. This means that feature selection and hyperparameter search remain internal to the training folds and that the final test error is evaluated only after the recursive search is complete. For the 261 scored non-bcc/fcc/hcp structures, the 80/20 split corresponds to 208 training and 53 test structures. The notebooks expose the starting dimensionalities of the descriptor spaces, including 1803 ACE features, 513 bond-specific BOP features, 419 or 287 SOAP features depending on the selected variant, and 32 structure-encoding features.

The primary regression metric reported in the archived outputs is the root-mean-square error,
\begin{equation}
\mathrm{RMSE}=\sqrt{\frac{1}{n}\sum_{i=1}^{n}\left(y_i^{\mathrm{pred}}-y_i^{\mathrm{DFT}}\right)^2},
\label{eq:rmse}
\end{equation}
with all values converted to meV/atom for discussion. The notebooks also evaluate the mean absolute error in parallel, but the human-readable stored outputs expose the RMSE tables most directly. Test-set predictions are inspected through parity plots, error histograms, and phase-resolved box plots.

\subsection{Finite-temperature thermodynamics}

To connect the predicted formation energies to experimentally accessible site occupancies, the repository performs Bragg--Williams calculations with a compound-energy-formalism representation for the complex phases. The thermodynamics notebook uses \texttt{pycef.cef\_minimization} to minimize the Gibbs free energy over sublattice occupancies at fixed temperature and composition. In the archived script this minimization is an end-member-only mean-field treatment: the phase energies are passed as sublattice end-member energies, no additional excess terms are exposed in the notebook text, and short-range order beyond Bragg--Williams is therefore neglected by construction. The archived production run is shown for \SI{1700}{\kelvin}, where the code evaluates the $R$, $P$, $M$, and C14-related free-energy surfaces over a composition grid. For the $R$ phase the sublattice model contains 11 site classes with multiplicities $[1,2,6,6,6,6,6,6,6,2,6]$, while $P$ and $M$ each use 12 sublattices and $\delta$ uses 14. For the $R$ phase the notebook prints the mixing contribution $\Delta G_{\mathrm{mix}}(x_{\mathrm{Fe}})$ explicitly; its most negative tabulated value is \SI{-53.372}{\milli\electronvolt\per\atom} at $x_{\mathrm{Fe}}=0.263158$. We emphasize that this number refers to the within-phase mixing contribution rather than, by itself, proving global stability against all competing phases. The resulting occupancies are compared with the experimental Fe--Mo site-fraction data stored in \texttt{Fe-Mo/ExperimentalData/FeMo.csv}, which span seven compositions between $x_{\mathrm{Mo}}=0.27958$ and 0.37978. % TODO: CITE [pycef or equivalent CEF implementation] % TODO: CITE [experimental R-phase occupancy measurements]
